%!TEX program = xelatex
%!TEX encoding = ISO-8859-1
\documentclass[12pt]{article}

\usepackage{amssymb}
\usepackage{amsmath,amsfonts,amsthm,amstext,mathabx}
\usepackage[brazil]{babel}
%\usepackage[latin1]{inputenc}
\usepackage{graphicx}
\graphicspath{{/home/jfreitas/Dropbox/imagens-latex/}{/Users/jfreitas/Dropbox/imagens-latex/}{D:/Dropbox/imagens-latex/}}
\usepackage{enumitem}
\usepackage{multicol}
\usepackage[all]{xy}

\setlength{\topmargin}{-1.0in}
\setlength{\oddsidemargin}{0in}
\setlength{\textheight}{10.1in}
\setlength{\textwidth}{6.5in}
\setlength{\baselineskip}{12mm}

\newcounter{exercicios}
\setcounter{exercicios}{0}
\newcommand{\questao}{
\addtocounter{exercicios}{1}
\noindent{\bf Exerc{\'\i}cio \arabic{exercicios}: }}

\newcommand{\equi}{\Leftrightarrow}
\newcommand{\bic}{\leftrightarrow}
\newcommand{\cond}{\rightarrow}
\newcommand{\impl}{\Rightarrow}
\newcommand{\nao}{\sim}
\newcommand{\sub}{\subseteq}
\newcommand{\e}{\ \wedge\ }
\newcommand{\ou}{\ \vee\ }
\newcommand{\vaz}{\emptyset}
\newcommand{\nsub}{\nsubset}
\renewcommand{\sin}{{\rm sen\,}}

\newcommand{\n}{\mathbb{N}}
\newcommand{\z}{\mathbb{Z}}
\newcommand{\real}{\mathbb{R}}
\newcommand{\vesp}{\vspace{0.2cm}}
\newcommand{\subne}{\subsetneqq}


\newcommand{\compcent}[1]{\vcenter{\hbox{$#1\circ$}}}
\newcommand{\comp}{\mathbin{\mathchoice
{\compcent\scriptstyle}{\compcent\scriptstyle}
{\compcent\scriptscriptstyle}{\compcent\scriptscriptstyle}}}

\begin{document}
\pagestyle{empty}

\begin{figure}[h]
        \begin{minipage}[c]{1.7cm}
        \includegraphics[width=1.7cm]{unb.pdf}
        \end{minipage}%
        \hspace{0pt}
        \begin{minipage}[c]{4in}
          {Universidade de Bras{\'\i}lia} \\
          {Departamento de Matem{\'a}tica}
\end{minipage}
\end{figure}
\vspace{-1cm}\hrule

\begin{center}
{\Large\bf {\'A}lgebra 1 - Turma D -- 2$^{o}$/2017} \\ \vspace{9pt} {\large\bf
  $1^{\underline{a}}$ Lista de Exerc{\'\i}cios -- Conjuntos}\\
\vspace{9pt} Prof. Jos{\'e} Ant{\^o}nio O. Freitas
\end{center}
\hrule

\vspace{.6cm}

\questao Dados os conjuntos $A = \{0,1,2\}$, $B = \{0,2,3\}$ e $C = \{0,1,2,3,4\}$, classifique as afirmações a seguir em verdadeira ou falsa:
\begin{multicols}{2}
	\begin{enumerate}[label={\alph*})]
		\item $A \sub B$
		\item $A \in C$
		\item $B \sub C$
		\item $\{0,2\} \sub B$
		\item $\{0,2\} \sub C$
		\item $1 \sub C$
		\item $\{1,4\} \in C$
		\item $\{0,3\} \sub B$
	\end{enumerate}	
\end{multicols}

\questao Dê exemplos de conjuntos não vazios $A$, $B$ e $C$ tais que:
\begin{enumerate}[label={\alph*})]
	\item $A \sub B$, $C \sub B$ e $A \cap C = \empty$.
	\item $A \sub B$, $C \nsub B$ e $A \cap C = \empty$.
	\item $A \sub C$, $A \ne C$ e $B \cap C = \empty$.
	\item $A \sub (B \cap C)$, $B \sub C$, $B \ne C$ e $A \ne C$.
\end{enumerate}

\vspace{.3cm}

\questao Em cada um dos seguintes itens, determine se a afirma{\c c}{\~a}o {\'e}
verdadeira ou falsa. Se for verdadeira, demonstre-a. Se for falsa, exiba um
exemplo mostrando que a afirma{\c c}{\~a}o {\'e} falsa.
\begin{enumerate}[label={\alph*})]
\item Se $x \in A$ e $A \sub B$, ent{\~a}o $x \in B$.

\item Se $A \nsub B$ e $B \subset C$, ent{\~a}o $A \nsub C$.

\item Se $A \nsub B$ e $B \nsub C$, ent{\~a}o $A \nsub C$.

\item Se $x \in A$ e $A \nsub B$, ent{\~a}o $x \notin B$.

\item Se $A \subset B$ e $x \notin B$, ent{\~a}o $x \notin A$.
\end{enumerate}

\vspace{.3cm}

\questao Determinar os elementos dos conjuntos $A$, $B$ e $E$ tais que:
\[
	A \cap B = \{b, c\}, \quad C_E(A) = \{d, e, f\} \quad \mbox{e}\quad C_E(B) = \{a, e, f\}.
\]

\vspace{.3cm}

\questao Dados os conjuntos $E = \{1,2,3,a,b,c,d\}$, $A = \{1,2,4,d\}$ e $B = \{a,2,4,b,5\}$, determine:
\begin{multicols}{2}
	\begin{enumerate}[label={\alph*})]
		\item $A \cup B$
		\item $A \cup B^C$
		\item $A^C \cup B$
		\item $(A \cup B)^C$
	\end{enumerate}
\end{multicols}

\vspace{.3cm}

\questao Dados os conjuntos $A = \{3,6,9,12,15,18\}$, $B = \{2,3,5,7,11,13,17,19\}$ e $C = \{1,2,3,4,5,\dots,20\}$ calcule:
\begin{multicols}{3}
	\begin{enumerate}[label={\alph*})]
		\item $A - B$
		\item $A - C$
		\item $B - C$
		\item $(A - B) - C$
		\item $A - (B - C)$
		\item $(A \cup B) - C$
		\item $A - (B \cap C)$
		\item $(A \cup C) - B$
		\item $(B \cap C) - A$
	\end{enumerate}
\end{multicols}

\vspace{.3cm}

\questao Demonstre que:
\begin{enumerate}[label={\alph*})]
\item Se $A \sub B$ e $C \sub D$, ent{\~a}o $A \cap C \sub B \cap D$.

\item Se $A \sub B$ e $C = B - A$, ent{\~a}o $A = B - C$.

\item Se $A \cap B = \vaz$ e $A \cup B = C$, ent{\~a}o $A = C - B$.

\item Se $A\cup B = E$, ent{\~a}o $C_E(A) \sub B$ (Aqui suponha que $A \sub E$ e $B \sub E$).

\item Se $A \cap B = \vaz$, ent{\~a}o $A \cup C_E(B) = C_E(B)$ (Aqui suponha que $A \sub E$ e $B \sub E$).
  
\item $A = B$ se, e somente se, $A - B = B - A$;

\item $A \sub B$ se, e somente se, $A - B = \vaz$;

\item $A \cup B = \vaz$ se, e somente se, $A = \vaz$ e $B = \vaz$;

\item $A \cup B = A \cap B$ se, e somente se,  $A = B$;

\item $C_E(A) \sub C_E(B)$ se, e somente se, $A \cup B = A$ (Aqui suponha que $A \sub E$ e $B \sub E$);

\item $C_E(A) \sub C_E(B)$ se, e somente se, $A \cap B = B$ (Aqui suponha que $A \sub E$ e $B \sub E$);

\item $A \sub B$ se, e somente se,  $A \cap B = A$.
\end{enumerate}

\vspace{.3cm}

\questao Demonstre as seguintes igualdades.
\begin{enumerate}[label={\alph*})]
\item $A \cup (C_E(A) \cap B) = A \cup B$ (Aqui suponha que $A \sub E$);

\item $A \cap (C_E(A) \cup B) = A \cap B$ (Aqui suponha que $A \sub E$);

\item $(A - B) - C = A - (B \cup C)$;

\item $A \cup (B - C) = (A \cup B) - (C - A)$;

\item $A \cap (B - C) = (A \cap B) - (A \cap C)$;

\item $A - (B \cup C) = (A - B) \cap (A - C)$;

\item $A - (B \cap C) = (A - B) \cup (A - C)$;

\item $(A \cup B) - C = (A - C) \cup (B - C)$;

\item $(A \cap B) - C = (A - C) \cap (B - C)$;

\item $(A \cap B) \cap (A - B) = (A - B) \cap (B - A) = \vaz$;

\item $(A - B) \cup (B - A) = (A \cup B) - (A \cap B)$;

\item $A - (A - B) = A \cap B$;

\item $(A - B) - B = A - B$;

\item $A \cup (B \cap (A \cup C)) = A \cup (B \cap C)$;

\item $C_E( A \cap (B \cup C)) = C_E(A) \cup (C_E(B)\cap C_E(C))$ (Aqui suponha que $A \sub E$, $B \sub E$ e $C \sub E$).

\item $C_E( A \cap B \cap C) = C_E(A) \cup C_E(B) \cup C_E(C)$ (Aqui suponha que $A \sub E$, $B \sub E$ e $C \sub E$).

\item $C_E( A \cup B \cup C) = C_E(A) \cap C_E(B) \cap C_E(C)$ (Aqui suponha que $A \sub E$, $B \sub E$ e $C \sub E$).
\end{enumerate}

\vspace{.3cm}

\questao Sejam $A$, $B$, $C$ e $D$ conjuntos.
\begin{enumerate}[label={\alph*})]
	\item Mostre que $A$ e $B$ são disjuntos se, e somente se, para todo conjunto não vazio $C$, $A \times C$ e $B \times C$ são disjuntos.
	\item Suponha $A \ne \emptyset$ e $C \ne \emptyset$, com $A \ne C$. Mostre que $A \sub B$ e $C \sub D$ se, e somente se, $A \times C \sub B \times D$.
\end{enumerate}

\vspace{.3cm}

\questao Sejam $X_1$, $X_1$, $Y_1$ e $Y_2$ subconjuntos contidos num conjunto $E$. Suponha que $X_1 \cup X_2 = E$, $Y_1 \cap Y_2 = \emptyset$, $X_1 \sub Y_1$ e $X_2 \sub Y_2$. Prove que $X_1 = Y_1$ e $X_2 = Y_2$.

\end{document}